\documentclass{article}
\usepackage{texmaths-symbols}
\begin{document}
% PROGRAMME 1
\noindent\textbf{Programme de construction 1}\par\smallskip\noindent
\begin{itemize}
    \item Tracer le segment $[\mathspoint{AB}]$ tel que
    $\mathspoint{AB}=10\text{ cm}$.
    \item Placer $\mathspoint{C}$ au milieu de $[\mathspoint{AB}]$.
    \item Tracer la droite $(\mathspoint{d})$ perpendiculaire à
    $[\mathspoint{AB}]$ passant par $\mathspoint{C}$.
    \item Sur la droite $(\mathspoint{d})$, placer les points
    $\mathspoint{D}$ et $\mathspoint{E}$ tels que
    $\mathspoint{DE}=6\text{ cm}$ et que $\mathspoint{C}$ soit le
    milieu de $[\mathspoint{DE}]$.
    \item Tracer la figure $\mathspoint{ADBE}$.
\end{itemize}
% PROGRAMME 2
\noindent\textbf{Programme de construction 2}\par\smallskip\noindent
\begin{itemize}
    \item Placer le point $\mathspoint{A}$.
    \item Tracer le cercle $\mathcal{C}$ de centre $\mathspoint{A}$ et de
    rayon $4\text{ cm}$.
    \item Tracer la droite $(\mathspoint{d})$ passant par $\mathspoint{A}$.
    \item Placer les points $\mathspoint{B}$ et $\mathspoint{C}$ tels que
    $\mathspoint{BC}=4\text{ cm}$ et que $\mathspoint{A}$ soit le milieu
    de $[\mathspoint{BC}]$.
    \item Tracer le cercle $\mathcal{D}$ de centre $\mathspoint{C}$ et de
    rayon $2\text{ cm}$.
    \item Tracer le cercle $\mathcal{E}$ de centre $\mathspoint{B}$ et de
    rayon $2\text{ cm}$.
\end{itemize}
% PROGRAMME 3
\noindent\textbf{Programme de construction 3}\par\smallskip\noindent
\begin{itemize}
    \item Tracer un rectangle $\mathspoint{ABCD}$ de largeur
    $\mathspoint{AB}=3\text{ cm}$ et de longueur $\mathspoint{BC}=4\text{ cm}$.
    \item Placer $\mathspoint{I}$, point d'intersection des diagonales.
    \item Tracer le cercle $\mathcal{C}$ de centre $\mathspoint{I}$ et de
    rayon $[\mathspoint{IA}]$.
    \item Que constates-tu ?
\end{itemize}
% PROGRAMME 4
\noindent\textbf{Programme de construction 4}\par\smallskip\noindent
\begin{itemize}
    \item Tracer le carré $\mathspoint{ABCD}$ de côté $6\text{ cm}$.
    \item Les diagonales se coupent en $\mathspoint{E}$.
    \item Tracer la droite $(\mathspoint{d})$ passant par $\mathspoint{E}$,
    perpendiculaire à $[\mathspoint{CD}]$.
    \item La droite $(\mathspoint{d})$ coupe $[\mathspoint{CD}]$ en
    $\mathspoint{F}$ et $[\mathspoint{BA}]$ en $\mathspoint{G}$.
    \item Placer le point $\mathspoint{H}$ sachant que $\mathspoint{F}$ est
    le milieu de $[\mathspoint{EH}]$.
    \item Placer le point $\mathspoint{I}$ sachant que $\mathspoint{G}$ est
    le milieu de $[\mathspoint{IE}]$.
    \item Tracer les figures $\mathspoint{ECHD}$ et $\mathspoint{IBEA}$,
    puis les colorier en rouge.
    \item Colorier les figures $\mathspoint{BCE}$ et $\mathspoint{AED}$ en
    vert.
    \item Quel est le nom des figures rouges ? Quel est le nom des figures
    vertes ?
\end{itemize}
% PROGRAMME 5
\noindent\textbf{Programme de construction 5}\par\smallskip\noindent
\begin{itemize}
    \item Tracer le segment $[\mathspoint{AB}]$ tel que
    $\mathspoint{AB}=7\text{ cm}$.
    \item Tracer la droite $(\mathspoint{d})$ perpendiculaire à
    $[\mathspoint{AB}]$ passant par le milieu de $[\mathspoint{AB}]$.
    \item La droite $(\mathspoint{d})$ coupe $[\mathspoint{AB}]$ en
    $\mathspoint{C}$.
    \item Placer le point $\mathspoint{E}$ tel que $\mathspoint{EC}=6\text{ cm}$.
    \item Tracer les segments $[\mathspoint{AE}]$ et $[\mathspoint{EB}]$.
    \item Tracer la droite $(\mathspoint{f})$ passant par $\mathspoint{E}$,
    perpendiculaire à $(\mathspoint{d})$.
    \item Tracer le cercle $\mathcal{C}$ de centre $\mathspoint{E}$ et de
    rayon $3{,}5\text{ cm}$.
    \item Le cercle $\mathcal{C}$ coupe la droite $(\mathspoint{f})$ en
    $\mathspoint{G}$ et $\mathspoint{H}$.
    \item Tracer les segments $[\mathspoint{AG}]$ et $[\mathspoint{BH}]$.
\end{itemize}
% PROGRAMME 6
\noindent\textbf{Programme de construction 6}\par\smallskip\noindent
\begin{itemize}
    \item Tracer une droite $(\mathspoint{d})$ et placer
    $\mathspoint{A}$ sur $(\mathspoint{d})$.
    \item Tracer le cercle $\mathcal{C}$ de centre $\mathspoint{A}$ et de
    rayon $4\text{ cm}$ ; il coupe $(\mathspoint{d})$ en
    $\mathspoint{B}$ et $\mathspoint{C}$.
    \item Tracer les demi-cercles $\mathcal{D}$, $\mathcal{E}$ et
    $\mathcal{F}$, respectivement de centre $\mathspoint{B}$,
    $\mathspoint{A}$ et $\mathspoint{C}$, et de rayon $2\text{ cm}$.
\end{itemize}
% PROGRAMME 7
\noindent\textbf{Programme de construction 7}\par\smallskip\noindent
\begin{itemize}
    \item Tracer le cercle $\mathcal{C}$ de centre $\mathspoint{A}$ et de
    rayon $4\text{ cm}$ (conserver l'écartement du compas).
    \item Placer le point $\mathspoint{B}$ sur le cercle. À l'aide du
    compas, marquer le point $\mathspoint{C}$ à partir de
    $\mathspoint{B}$, le point $\mathspoint{D}$ à partir de
    $\mathspoint{C}$, le point $\mathspoint{E}$ à partir de
    $\mathspoint{D}$, le point $\mathspoint{F}$ à partir de
    $\mathspoint{E}$ et le point $\mathspoint{G}$ à partir de
    $\mathspoint{F}$.
    \item Tracer les triangles $\mathspoint{EGC}$ et $\mathspoint{FBD}$.
    \item Nommer l'hexagone obtenu $\mathspoint{HIJKLM}$.
    \item Tracer les triangles $\mathspoint{LHJ}$ et $\mathspoint{MIK}$.
\end{itemize}
% PROGRAMME 8
\noindent\textbf{Programme de construction 8}\par\smallskip\noindent
\begin{itemize}
    \item Tracer le rectangle $\mathspoint{ABCD}$ de longueur
    $\mathspoint{BC}=5\text{ cm}$ et de largeur $\mathspoint{AB}=2\text{ cm}$.
    \item Tracer la droite $(\mathspoint{d})$ passant par
    $\mathspoint{A}$ et $\mathspoint{C}$.
    \item Placer $\mathspoint{E}$ sur la droite $(\mathspoint{d})$ sachant
    que $[\mathspoint{AC}]=[\mathspoint{CE}]$.
    \item Placer $\mathspoint{F}$, milieu de $[\mathspoint{CE}]$.
    \item Tracer le cercle $\mathcal{C}$ de centre $\mathspoint{F}$ et de
    rayon $[\mathspoint{CF}]$.
\end{itemize}
% PROGRAMME 9
\noindent\textbf{Programme de construction 9}\par\smallskip\noindent
\begin{itemize}
    \item Tracer le cercle $\mathcal{C}$ de centre $\mathspoint{A}$ et de
    rayon $4\text{ cm}$.
    \item Placer le point $\mathspoint{B}$ sur le cercle. À l'aide du
    compas, marquer le point $\mathspoint{C}$ à partir de
    $\mathspoint{B}$, le point $\mathspoint{D}$ à partir de
    $\mathspoint{C}$, le point $\mathspoint{E}$ à partir de
    $\mathspoint{D}$, le point $\mathspoint{F}$ à partir de
    $\mathspoint{E}$ et le point $\mathspoint{G}$ à partir de
    $\mathspoint{F}$.
    \item Placer la pointe du compas sur le point $\mathspoint{B}$ et
    relier les deux points adjacents à $\mathspoint{B}$. Faire de même pour
    chaque point.
\end{itemize}
% PROGRAMME 10
\noindent\textbf{Programme de construction 10}\par\smallskip\noindent
\begin{itemize}
    \item Tracer une droite $(\mathspoint{d})$ et placer le point
    $\mathspoint{A}$ dessus.
    \item Tracer sur la droite $(\mathspoint{d})$ un arc de cercle
    $\mathcal{C}$ de centre $\mathspoint{A}$ et de rayon $3\text{ cm}$. Il
    coupe la droite $(\mathspoint{d})$ en $\mathspoint{B}$ et
    $\mathspoint{C}$.
    \item Tracer le segment $[\mathspoint{AD}]$ à l'opposé de
    $\mathspoint{C}$, sachant que $\mathspoint{AD}=3{,}5\text{ cm}$ et que
    $[\mathspoint{AD}]$ est perpendiculaire à $(\mathspoint{d})$.
    \item Tracer $[\mathspoint{BD}]$ et $[\mathspoint{DC}]$.
\end{itemize}
% PROGRAMME 11
\noindent\textbf{Programme de construction 11}\par\smallskip\noindent
\begin{itemize}
    \item Tracer un carré $\mathspoint{ABCD}$ de côté $3\text{ cm}$.
    \item Les diagonales se croisent en $\mathspoint{E}$.
    \item Tracer le cercle $\mathcal{A}$ de centre $\mathspoint{A}$ et de
    rayon $[\mathspoint{AE}]$.
    \item Tracer le cercle $\mathcal{B}$ de centre $\mathspoint{B}$ et de
    rayon $[\mathspoint{BE}]$.
    \item Tracer le cercle $\mathcal{C}$ de centre $\mathspoint{C}$ et de
    rayon $[\mathspoint{CE}]$.
    \item Tracer le cercle $\mathcal{D}$ de centre $\mathspoint{D}$ et de
    rayon $[\mathspoint{DE}]$.
\end{itemize}
% PROGRAMME 12
\noindent\textbf{Programme de construction 12}\par\smallskip\noindent
\begin{itemize}
    \item Tracer le segment $[\mathspoint{AB}]$ tel que
    $\mathspoint{AB}=8\text{ cm}$.
    \item Placer $\mathspoint{C}$ au milieu de $[\mathspoint{AB}]$.
    \item Tracer la droite $(\mathspoint{d})$ passant par
    $\mathspoint{C}$, perpendiculaire à $[\mathspoint{AB}]$.
    \item Placer $\mathspoint{D}$ et $\mathspoint{E}$ sur
    $(\mathspoint{d})$ sachant que $\mathspoint{C}$ est le milieu de
    $[\mathspoint{DE}]$ et que $\mathspoint{DE}=4\text{ cm}$.
    \item Placer $\mathspoint{F}$ et $\mathspoint{G}$ tels que
    $\mathspoint{FG}=4\text{ cm}$ et que $\mathspoint{C}$ soit le milieu
    de $[\mathspoint{FG}]$.
    \item Tracer le carré $\mathspoint{FDGE}$ et le losange
    $\mathspoint{DAEB}$.
\end{itemize}
% PROGRAMME 13
\noindent\textbf{Programme de construction 13}\par\smallskip\noindent
\begin{itemize}
    \item Tracer un carré $\mathspoint{ABCD}$ de côté $5\text{ cm}$.
    \item Placer le point $\mathspoint{E}$, milieu de $[\mathspoint{AB}]$.
    \item Tracer le segment $[\mathspoint{EF}]$ de longueur $3\text{ cm}$,
    perpendiculaire à $[\mathspoint{AB}]$ et à l'extérieur du carré.
    \item Tracer le triangle $\mathspoint{AFB}$.
    \item Placer les points $\mathspoint{G}$ et $\mathspoint{H}$ tels que
    $\mathspoint{DG}=2\text{ cm}$, $\mathspoint{GH}=1\text{ cm}$ et
    $\mathspoint{HC}=2\text{ cm}$.
    \item Tracer le carré $\mathspoint{GIJH}$ à l'intérieur de
    $\mathspoint{ABCD}$.
\end{itemize}
% PROGRAMME 14
\noindent\textbf{Programme de construction 14}\par\smallskip\noindent
\begin{itemize}
    \item Tracer une droite $(\mathspoint{d})$ et placer un point
    $\mathspoint{A}$ sur $(\mathspoint{d})$.
    \item Placer $\mathspoint{B}$ sachant que $\mathspoint{AB}=2\text{ cm}$.
    \item Tracer les cercles $\mathcal{C}$ de centre $\mathspoint{A}$ et
    $\mathcal{D}$ de centre $\mathspoint{B}$, de rayon $1{,}5\text{ cm}$.
    \item Les cercles $\mathcal{C}$ et $\mathcal{D}$ se croisent en
    $\mathspoint{C}$ et $\mathspoint{D}$.
    \item Tracer la droite $(\mathspoint{e})$ passant par
    $\mathspoint{C}$ et $\mathspoint{D}$.
    \item Tracer le quadrilatère $\mathspoint{ACBD}$.
\end{itemize}
% PROGRAMME 15
\noindent\textbf{Programme de construction 15}\par\smallskip\noindent
\begin{itemize}
    \item Tracer le segment $[\mathspoint{AB}]$ tel que
    $\mathspoint{AB}=10\text{ cm}$.
    \item Placer $\mathspoint{C}$, milieu de $[\mathspoint{AB}]$.
    \item Placer $\mathspoint{D}$, milieu de $[\mathspoint{AC}]$.
    \item Placer $\mathspoint{E}$, milieu de $[\mathspoint{CB}]$.
    \item Tracer la droite $(\mathspoint{d})$ passant par
    $\mathspoint{C}$, perpendiculaire à $[\mathspoint{AB}]$.
    \item Placer $\mathspoint{F}$ sachant que $\mathspoint{CF}=5\text{ cm}$.
    \item Tracer les triangles $\mathspoint{DFE}$ et $\mathspoint{AFB}$.
    \item Tracer la droite $(\mathspoint{e})$ passant par
    $\mathspoint{F}$, parallèle à $[\mathspoint{AB}]$.
    \item Tracer la droite $(\mathspoint{f})$ passant par
    $\mathspoint{A}$, perpendiculaire à $(\mathspoint{e})$ ; elle coupe
    $(\mathspoint{e})$ en $\mathspoint{G}$.
    \item Tracer la droite $(\mathspoint{g})$ passant par
    $\mathspoint{B}$, perpendiculaire à $(\mathspoint{e})$ ; elle coupe
    $(\mathspoint{e})$ en $\mathspoint{H}$.
\end{itemize}
% PROGRAMME 16
\noindent\textbf{Programme de construction 16}\par\smallskip\noindent
\begin{itemize}
    \item Tracer un rectangle $\mathspoint{ABCD}$ de largeur
    $\mathspoint{AB}=4\text{ cm}$ et de longueur $\mathspoint{BC}=5\text{ cm}$.
    \item Les diagonales se coupent en $\mathspoint{E}$.
    \item Tracer le cercle $\mathcal{C}$ de centre $\mathspoint{E}$ et de
    rayon $[\mathspoint{EA}]$.
    \item Tracer le cercle $\mathcal{D}$ de centre $\mathspoint{E}$ et de
    rayon $2\text{ cm}$.
\end{itemize}
% PROGRAMME 17
\noindent\textbf{Programme de construction 17}\par\smallskip\noindent
\begin{itemize}
    \item Tracer les droites $(\mathspoint{d})$ et $(\mathspoint{e})$
    perpendiculaires entre elles. Elles se coupent en $\mathspoint{A}$.
    \item Tracer le segment $[\mathspoint{BC}]$ tel que
    $\mathspoint{BC}=4\text{ cm}$ et que $\mathspoint{A}$ soit le milieu
    de $[\mathspoint{BC}]$.
    \item Tracer les demi-cercles $\mathcal{C}$ et $\mathcal{D}$,
    respectivement de centre $\mathspoint{B}$ et de centre
    $\mathspoint{C}$, de rayon $2\text{ cm}$.
    \item Le cercle $\mathcal{C}$ coupe la droite $(\mathspoint{d})$ en
    $\mathspoint{E}$ et le cercle $\mathcal{D}$ coupe $(\mathspoint{d})$
    en $\mathspoint{F}$.
    \item À l'opposé des demi-cercles, placer le point $\mathspoint{G}$ sur
    la droite $(\mathspoint{e})$ sachant que $\mathspoint{AG}=3\text{ cm}$.
    \item Tracer les segments $[\mathspoint{FG}]$ et $[\mathspoint{DG}]$.
\end{itemize}
% PROGRAMME 18
\noindent\textbf{Programme de construction 18}\par\smallskip\noindent
\begin{itemize}
    \item Tracer un rectangle $\mathspoint{ABCD}$ de longueur
    $\mathspoint{AB}=7\text{ cm}$ et de largeur $\mathspoint{BC}=2\text{ cm}$.
    \item Placer $\mathspoint{E}$, milieu de $[\mathspoint{AB}]$.
    \item Tracer le cercle $\mathcal{C}$ de centre $\mathspoint{E}$ et de
    rayon $4\text{ cm}$.
    \item Tracer la droite $(\mathspoint{d})$ passant par
    $\mathspoint{E}$, perpendiculaire à $[\mathspoint{AB}]$.
    \item Tracer les diamètres du cercle $\mathcal{C}$ passant par
    $\mathspoint{C}$ et $\mathspoint{D}$.
\end{itemize}
% PROGRAMME 19
\noindent\textbf{Programme de construction 19}\par\smallskip\noindent
\begin{itemize}
    \item Tracer une droite $(\mathspoint{d})$, puis une droite
    $(\mathspoint{e})$ parallèle à $(\mathspoint{d})$ et située au-dessus,
    à $3\text{ cm}$.
    \item Placer le point $\mathspoint{A}$ sur la droite $(\mathspoint{e})$.
    \item Tracer le cercle $\mathcal{C}$ de centre $\mathspoint{A}$ et de
    rayon $3{,}5\text{ cm}$.
    \item Le cercle $\mathcal{C}$ coupe la droite $(\mathspoint{e})$ en
    $\mathspoint{B}$ et $\mathspoint{C}$, et la droite
    $(\mathspoint{d})$ en $\mathspoint{D}$ et $\mathspoint{E}$.
    \item Tracer la figure $\mathspoint{BCDE}$.
    \item Quel est le nom de cette figure ?
\end{itemize}
% PROGRAMME 20
\noindent\textbf{Programme de construction 20}\par\smallskip\noindent
\begin{itemize}
    \item Tracer un carré $\mathspoint{ABCD}$ de côté $3\text{ cm}$.
    \item Tracer le cercle $\mathcal{A}$ de centre $\mathspoint{A}$ et de
    rayon $1{,}5\text{ cm}$.
    \item Tracer le cercle $\mathcal{B}$ de centre $\mathspoint{B}$ et de
    rayon $1{,}5\text{ cm}$.
    \item Tracer le cercle $\mathcal{C}$ de centre $\mathspoint{C}$ et de
    rayon $1{,}5\text{ cm}$.
    \item Tracer le cercle $\mathcal{D}$ de centre $\mathspoint{D}$ et de
    rayon $1{,}5\text{ cm}$.
\end{itemize}
% PROGRAMME 21
\noindent\textbf{Programme de construction 21}\par\smallskip\noindent
\begin{itemize}
    \item Tracer le cercle $\mathcal{C}$ de centre $\mathspoint{A}$ et de
    rayon $4\text{ cm}$.
    \item Tracer les diamètres $[\mathspoint{EB}]$ et $[\mathspoint{CD}]$.
    \item Le segment $[\mathspoint{EB}]$ est perpendiculaire au segment
    $[\mathspoint{CD}]$.
    \item Tracer le carré $\mathspoint{ADBC}$.
\end{itemize}
% PROGRAMME 22
\noindent\textbf{Programme de construction 22}\par\smallskip\noindent
\begin{itemize}
    \item Tracer la droite $(\mathspoint{d})$.
    \item Placer $\mathspoint{A}$ et $\mathspoint{B}$ sur
    $(\mathspoint{d})$ tels que $\mathspoint{AB}=4\text{ cm}$.
    \item Tracer le cercle $\mathcal{A}$ de centre $\mathspoint{A}$ et de
    rayon $2\text{ cm}$.
    \item Tracer le cercle $\mathcal{B}$ de centre $\mathspoint{B}$ et de
    rayon $2\text{ cm}$.
    \item Placer $\mathspoint{C}$, milieu de $[\mathspoint{AB}]$.
    \item Tracer la droite $(\mathspoint{e})$ perpendiculaire à
    $(\mathspoint{d})$ passant par $\mathspoint{C}$.
    \item Placer les points $\mathspoint{D}$ et $\mathspoint{E}$ sachant
    que $\mathspoint{DC}=2\text{ cm}$ et $\mathspoint{CE}=3\text{ cm}$.
    \item Tracer le quadrilatère $\mathspoint{EBDA}$.
\end{itemize}
% PROGRAMME 23
\noindent\textbf{Programme de construction 23}\par\smallskip\noindent
\begin{itemize}
    \item Tracer un carré $\mathspoint{ABCD}$ de côté $6\text{ cm}$.
    \item Sur $[\mathspoint{AB}]$, placer les points $\mathspoint{E}$ et
    $\mathspoint{F}$ tels que $\mathspoint{AE}=\mathspoint{FB}=2\text{ cm}$.
    \item Sur $[\mathspoint{BC}]$, placer les points $\mathspoint{G}$ et
    $\mathspoint{H}$ tels que $\mathspoint{BG}=\mathspoint{HC}=2\text{ cm}$.
    \item Sur $[\mathspoint{CD}]$, placer les points $\mathspoint{I}$ et
    $\mathspoint{J}$ tels que $\mathspoint{CI}=\mathspoint{JD}=2\text{ cm}$.
    \item Sur $[\mathspoint{DA}]$, placer les points $\mathspoint{K}$ et
    $\mathspoint{L}$ tels que $\mathspoint{DK}=\mathspoint{LA}=2\text{ cm}$.
    \item Tracer les segments $[\mathspoint{EJ}]$, $[\mathspoint{FI}]$,
    $[\mathspoint{LG}]$ et $[\mathspoint{KH}]$.
    \item Colorier en orange les rectangles $\mathspoint{EFIJ}$ et
    $\mathspoint{LGHK}$.
    \item Colorier en violet les carrés.
\end{itemize}
% PROGRAMME 24
\noindent\textbf{Programme de construction 24}\par\smallskip\noindent
\begin{itemize}
    \item Tracer le segment $[\mathspoint{AB}]$ tel que
    $\mathspoint{AB}=6\text{ cm}$.
    \item Tracer le triangle équilatéral $\mathspoint{ABC}$.
    \item Placer $\mathspoint{D}$, milieu de $[\mathspoint{AB}]$.
    \item Tracer la droite $(\mathspoint{d})$ passant par
    $\mathspoint{C}$ et $\mathspoint{D}$.
    \item Placer le point $\mathspoint{E}$, milieu de $[\mathspoint{DC}]$.
    \item Tracer le cercle $\mathcal{C}$ de centre $\mathspoint{E}$ et de
    rayon $1{,}3\text{ cm}$.
\end{itemize}
% PROGRAMME 25
\noindent\textbf{Programme de construction 25}\par\smallskip\noindent
\begin{itemize}
    \item Tracer un rectangle $\mathspoint{ABCD}$ de longueur
    $\mathspoint{AB}=6\text{ cm}$ et de largeur $\mathspoint{BC}=2\text{ cm}$.
    \item Placer $\mathspoint{E}$, milieu de $[\mathspoint{AD}]$, et
    $\mathspoint{F}$, milieu de $[\mathspoint{BC}]$.
    \item À droite de $[\mathspoint{BC}]$, tracer le segment
    $[\mathspoint{FJ}]$ tel que $\mathspoint{FJ}=2\text{ cm}$,
    perpendiculaire à $[\mathspoint{BC}]$.
    \item À gauche de $[\mathspoint{AD}]$, tracer le segment
    $[\mathspoint{IE}]$ tel que $\mathspoint{IE}=2\text{ cm}$,
    perpendiculaire à $[\mathspoint{AD}]$.
    \item Placer le point $\mathspoint{G}$ sur $[\mathspoint{DC}]$ à
    $2\text{ cm}$ de $\mathspoint{D}$.
    \item Placer le point $\mathspoint{H}$ sur $[\mathspoint{DC}]$ à
    $2\text{ cm}$ de $\mathspoint{C}$.
    \item En dessous de $[\mathspoint{DC}]$, tracer $[\mathspoint{HL}]$ tel
    que $\mathspoint{HL}=1\text{ cm}$, perpendiculaire à $[\mathspoint{DC}]$.
    \item En dessous de $[\mathspoint{DC}]$, tracer $[\mathspoint{GK}]$ tel
    que $\mathspoint{GK}=1\text{ cm}$, perpendiculaire à $[\mathspoint{DC}]$.
    \item Tracer le segment $[\mathspoint{NK}]$ tel que
    $\mathspoint{NK}=4\text{ cm}$, perpendiculaire à $[\mathspoint{GK}]$.
    \item Tracer le segment $[\mathspoint{LM}]$ tel que
    $\mathspoint{LM}=4\text{ cm}$, perpendiculaire à $[\mathspoint{HL}]$.
    \item Tracer la droite $(\mathspoint{d})$ passant par
    $\mathspoint{G}$ et $\mathspoint{H}$.
    \item Placer $\mathspoint{O}$ sur $(\mathspoint{d})$ sachant que
    $\mathspoint{OD}=1{,}5\text{ cm}$.
    \item Placer $\mathspoint{P}$ sur $(\mathspoint{d})$ sachant que
    $\mathspoint{CP}=1{,}5\text{ cm}$.
    \item Tracer les segments $[\mathspoint{JP}]$, $[\mathspoint{PM}]$,
    $[\mathspoint{LC}]$, $[\mathspoint{DK}]$, $[\mathspoint{IO}]$ et
    $[\mathspoint{ON}]$.
    \item Colorier les triangles en bleu foncé et les autres figures en
    bleu clair.
\end{itemize}
% PROGRAMME 26
\noindent\textbf{Programme de construction 26}\par\smallskip\noindent
\begin{itemize}
    \item Tracer un cercle de centre $\mathspoint{C}$ et de rayon
    $3\text{ cm}$ (conserver l'écartement du compas).
    \item Placer $\mathspoint{B}$ sur le cercle puis, à l'aide du compas,
    reporter la longueur sur le cercle à partir de $\mathspoint{B}$ pour
    placer $\mathspoint{H}$, à partir de $\mathspoint{H}$ pour
    $\mathspoint{D}$, à partir de $\mathspoint{D}$ pour
    $\mathspoint{E}$, à partir de $\mathspoint{E}$ pour
    $\mathspoint{F}$ et à partir de $\mathspoint{F}$ pour
    $\mathspoint{G}$.
    \item Tracer les triangles $\mathspoint{GHE}$, $\mathspoint{GBH}$,
    $\mathspoint{EHD}$ et $\mathspoint{GEF}$.
    \item Colorier le triangle équilatéral en violet.
    \item Colorier les triangles isocèles en orange.
\end{itemize}
% PROGRAMME 27
\noindent\textbf{Programme de construction 27}\par\smallskip\noindent
\begin{itemize}
    \item Tracer une droite $(\mathspoint{d})$.
    \item Placer les points $\mathspoint{A}$, $\mathspoint{B}$,
    $\mathspoint{C}$ et $\mathspoint{D}$ sur $(\mathspoint{d})$ tels que
    $\mathspoint{AB}=1\text{ cm}$, $\mathspoint{BC}=5\text{ cm}$ et
    $\mathspoint{CD}=1\text{ cm}$.
    \item Tracer les droites perpendiculaires à $(\mathspoint{d})$ :
    \begin{itemize}
        \item[(a)] passant par $\mathspoint{A}$ ;
        \item[(b)] passant par $\mathspoint{B}$ ;
        \item[(c)] passant par $\mathspoint{C}$ ;
        \item[(e)] passant par $\mathspoint{D}$.
    \end{itemize}
    \item Tracer la droite $(\mathspoint{f})$ en dessous de
    $(\mathspoint{d})$, parallèle à $(\mathspoint{d})$ et distante de
    $4\text{ cm}$.
    \item La droite $(\mathspoint{f})$ coupe la droite (a) en
    $\mathspoint{E}$, la droite (b) en $\mathspoint{F}$, la droite (c) en
    $\mathspoint{G}$ et la droite (e) en $\mathspoint{H}$.
    \item Tracer le quadrilatère $\mathspoint{BDGE}$.
    \item Quel est le nom de cette figure ?
\end{itemize}
% PROGRAMME 28
\noindent\textbf{Programme de construction 28}\par\smallskip\noindent
\begin{itemize}
    \item Tracer une droite $(\mathspoint{d})$ et une droite
    $(\mathspoint{e})$ parallèle, au-dessus de $(\mathspoint{d})$ et à
    $4\text{ cm}$.
    \item Placer les points $\mathspoint{A}$, $\mathspoint{B}$ et
    $\mathspoint{C}$ sur $(\mathspoint{d})$ tels que
    $\mathspoint{AB}=\mathspoint{BC}=4\text{ cm}$.
    \item Tracer les droites perpendiculaires à $(\mathspoint{d})$ :
    \begin{itemize}
        \item[(a)] passant par $\mathspoint{A}$ ;
        \item[(b)] passant par $\mathspoint{B}$ ;
        \item[(c)] passant par $\mathspoint{C}$.
    \end{itemize}
    \item La droite (b) coupe la droite $(\mathspoint{e})$ en
    $\mathspoint{D}$, la droite (a) coupe $(\mathspoint{e})$ en
    $\mathspoint{E}$ et la droite (c) coupe $(\mathspoint{e})$ en
    $\mathspoint{F}$.
    \item Sur la droite (a), placer les points $\mathspoint{G}$ et
    $\mathspoint{H}$ tels que $\mathspoint{GE}=3\text{ cm}$ et
    $\mathspoint{EH}=1\text{ cm}$.
    \item Sur la droite (c), placer les points $\mathspoint{I}$ et
    $\mathspoint{J}$ tels que $\mathspoint{IF}=3\text{ cm}$ et
    $\mathspoint{FJ}=1\text{ cm}$.
    \item Tracer les triangles $\mathspoint{GBA}$ et $\mathspoint{BIC}$.
    Les colorier en marron.
    \item Tracer les parallélogrammes $\mathspoint{GHDB}$ et
    $\mathspoint{BDJI}$. Les colorier en jaune.
\end{itemize}
% PROGRAMME 29
\noindent\textbf{Programme de construction 29}\par\smallskip\noindent
\begin{itemize}
    \item Tracer le cercle $\mathcal{C}$ de centre $\mathspoint{A}$ et de
    rayon $4\text{ cm}$.
    \item Placer un point $\mathspoint{B}$ sur le cercle $\mathcal{C}$.
    \item À l'aide du compas, tracer un arc de cercle de centre
    $\mathspoint{B}$ et de rayon $[\mathspoint{AB}]$. Il croise le cercle
    $\mathcal{C}$ en $\mathspoint{C}$ et $\mathspoint{D}$.
    \item Placer le point $\mathspoint{E}$, milieu de $[\mathspoint{AB}]$.
    \item Tracer le cercle $\mathcal{B}$ de centre $\mathspoint{E}$ et de
    rayon $2\text{ cm}$.
\end{itemize}
% PROGRAMME 30
\noindent\textbf{Programme de construction 30}\par\smallskip\noindent
\begin{itemize}
    \item Tracer une droite $(\mathspoint{d})$ et placer le point
    $\mathspoint{A}$ sur $(\mathspoint{d})$.
    \item Tracer la droite $(\mathspoint{e})$ passant par
    $\mathspoint{A}$, perpendiculaire à $(\mathspoint{d})$.
    \item Placer le point $\mathspoint{C}$ sur $(\mathspoint{e})$ tel que
    $\mathspoint{AC}=4\text{ cm}$.
    \item Placer le point $\mathspoint{B}$ sur $(\mathspoint{d})$ sachant
    que $\mathspoint{AB}=3\text{ cm}$, puis tracer $[\mathspoint{BC}]$.
    \item Tracer la droite $(\mathspoint{f})$ parallèle à
    $(\mathspoint{d})$ et passant par $\mathspoint{C}$.
    \item Placer le point $\mathspoint{D}$ sur $(\mathspoint{d})$ sachant
    que $\mathspoint{DA}=4\text{ cm}$.
    \item Placer le point $\mathspoint{E}$ sur $(\mathspoint{f})$ sachant
    que $\mathspoint{EC}=4\text{ cm}$.
    \item Placer $\mathspoint{F}$ sur $(\mathspoint{d})$ sachant que
    $\mathspoint{BF}=2\text{ cm}$.
    \item Placer $\mathspoint{G}$ sur $(\mathspoint{f})$ sachant que
    $\mathspoint{CG}=2\text{ cm}$.
    \item Placer $\mathspoint{H}$ sur $(\mathspoint{d})$ sachant que
    $\mathspoint{HD}=1\text{ cm}$.
    \item Placer $\mathspoint{I}$ sur $(\mathspoint{f})$ sachant que
    $\mathspoint{IE}=1\text{ cm}$.
    \item Colorier en respectant le code couleur suivant : triangle en
    vert, rectangle en bleu, carré en jaune, parallélogramme en rouge.
\end{itemize}
% PROGRAMME 31
\noindent\textbf{Programme de construction 31}\par\smallskip\noindent
\begin{itemize}
    \item Tracer un carré $\mathspoint{ABCD}$ de côté $5\text{ cm}$.
    \item Les diagonales se coupent en $\mathspoint{E}$.
    \item Placer $\mathspoint{F}$, milieu de $[\mathspoint{AB}]$.
    \item Tracer le cercle $\mathcal{C}$ de centre $\mathspoint{E}$ et de
    rayon $[\mathspoint{EF}]$.
\end{itemize}
% PROGRAMME 32
\noindent\textbf{Programme de construction 32}\par\smallskip\noindent
\begin{itemize}
    \item Tracer le segment $[\mathspoint{AD}]$ tel que
    $\mathspoint{AD}=10\text{ cm}$.
    \item Placer $\mathspoint{B}$ et $\mathspoint{C}$ sur
    $[\mathspoint{AD}]$ sachant que
    $\mathspoint{AB}=\mathspoint{CD}=2\text{ cm}$ et
    $\mathspoint{BC}=6\text{ cm}$.
    \item Tracer la droite $(\mathspoint{d})$ passant par
    $\mathspoint{B}$, perpendiculaire à $[\mathspoint{AD}]$.
    \item Tracer la droite $(\mathspoint{e})$ passant par
    $\mathspoint{C}$, perpendiculaire à $[\mathspoint{AD}]$.
    \item Tracer $\mathspoint{AEFD}$. Quel est le nom de cette figure ?
\end{itemize}
\end{document}